\documentclass[12pt, a4paper, twoside, english, brazil]{article}
\usepackage[top=2cm, bottom=2cm, left=2cm, right=2cm]{geometry}

\usepackage{mathtools}
\usepackage{amsfonts}
\usepackage{amssymb}
\usepackage{amsmath}
\usepackage[utf8]{inputenc}
\usepackage[portuguese]{babel}
\usepackage[T1]{fontenc}
\usepackage{gensymb}

\usepackage{enumerate}
\usepackage{multicol}
\usepackage{tikz}
\usepackage{pgfplots}

\newenvironment{solution}[1][]{\noindent\textbf{Solução#1:} }

% -----------------------------------------------------------------------------

\title{Matemática Discreta e Lógica Matemática - Lista 03}
\author{Davi Reis - Funções do Segundo Grau (24/04/2026)}
\date{}

\begin{document}

\maketitle

\begin{enumerate}[\bfseries 1.]
    % Questão 1
  \item Seja a função $\text{altura}(t)=-4.9t^2+50t+2$, que calcula a
    altura ($m$), de um objeto lançado verticalmente para cima, em
    função do tempo $t$ ($s$), responda:
  \begin{enumerate}[\bfseries a)]
    \item Qual a altura inicial do objeto?
    \item A altura em $t=2s$ e $t=4s$?
    \item Em que instante o objeto toca o solo?
    \item Qual foi a altura máxima atingida pelo objeto?
  \end{enumerate}

  \begin{solution}
  \begin{enumerate}[\bfseries a)]
    \item $\text{altura}(0)=0+0+2=2m$
    \item $\text{altura}(2)=-4.9\cdot4+50\cdot2+2=82.4m$\\
      $\text{altura}(4)=-4.9\cdot16+50\cdot4+2=123.6m$
    \item $-4.9t^2+50t+2=0\implies \Delta = 50^2-4\cdot(-4.9)\cdot2=2539.2$\\
      $t=\frac{-50\pm\sqrt{2539.2}}{2\cdot(-4.9)}\approx\frac{-50\pm50.39047529}{−9.8}\approx\frac{100.39}{9.8}\approx10.24s$
    \item $t_{max}=\frac{-50}{2\cdot(-4.9)}\approx5.10s$\\
      $\text{altura}_{max}=-4.9\cdot(5.10)^2+50\cdot5.10+2\approx129.55m$
  \end{enumerate}
\end{solution}

% Questão 2
\item O saldo da conta corrente (milhares) de Pedro, ao longo do ano
de 2026, é dado em função dos meses $m$ (meses: janeiro = 1,
fevereiro = 2 e etc).
$$\text{saldo}(m)=m^2-10m+16$$
Responda:
\begin{enumerate}[\bfseries a)]
\item Qual o saldo da conta em dezembro de 2025 ($m=0$)?
\item Quanto ele terá em dezembro de 2026?
\item Quando ele teve o saldo 0 (zero)?
\item Quanto foi o saldo mínimo?
\end{enumerate}

\begin{solution}
\begin{enumerate}[\bfseries a)]
\item $\text{saldo}(0)=16\text{ mil}$
\item $\text{saldo}(12)=12^2-10\cdot12+16=40\text{ mil}$
\item $m^2-10m+16=0\implies \Delta=(-10)^2-4\cdot1\cdot16=36$\\
  $m=\frac{-(-10)\pm 6}{2}\implies \text{fevereiro (2) e agosto (8)}
\item $m_{min}=\frac{-(-10)}{2\cdot 1}=5\text{ (maio)}$\\
  $\text{saldo}_{min}=5^2-10\cdot 5+16=-9\text{ mil}$
\end{enumerate}
\end{solution}

% Questão 3
\item Determinado satélite demora 10h para completar sua orbita em
torno da terra. Durante certo período o satélite fica encoberto pela
Terra, não recebendo luz solar, o que faz sua temperatura despencar.
A temperatura média $T$ ($C$) de seus instrumentos é descrita pelo
modelo $T(h)=2h^2-20h+32$ em função das horas $h$. Responda:
\begin{enumerate}[\bfseries a)]
\item Calcule a temperatura nas horas $h=1$, $h=7$ e $h=9$.
\item Qual a temperatura mínima?
\item Quando a temperatura fica igual a 0 (zero)?
\item Para quais intervalos de tempo $h$, a temperatura é positiva
($T(h)>0$) e quando é negativa ($T(h)<0$). Esboce o gráfico.
\end{enumerate}

\begin{solution}
\begin{enumerate}[\bfseries a)]
\item $T(1)=2\cdot 1^2-20\cdot1+32=14\degree\text{C}$\\
$T(7)=2\cdot7^2-20\cdot 7+32=-10\degree\text{C}$\\
$T(9)=2\cdot9^2-20\cdot 9+32=14\degree\text{C}$
\item $h_{min}=\frac{-(-20)}{2\cdot 2}=5h$\\
$T_{min}=2\cdot(5)^2-20\cdot 5+32=-18\degree\text{C}$
\item $2h^2-20h+32=0\implies h^2-10h+16=0\implies h=
\begin{cases}2h\\8h
\end{cases}$
\item Fica positiva para $h<2$ ou para $h>8$, e fica negativa para $2<h<8$:
\begin{tikzpicture}
\begin{axis}[
    axis lines = center,
    xlabel = \(x\),ylabel ={\(y\)},
    xtick distance=2,
    grid,
  ]

  \addplot [domain=-10:20,color=blue]{2 * x ^ 2 - 20 * x + 32};

  \addlegendentry{$T(h)$}
\end{axis}
\end{tikzpicture}
\end{enumerate}
\end{solution}

\end{enumerate}
\end{document}
